\section{把三类评估拼成完整能力画像}

\videofigure{figures/long-horizon-implications.jpg}{METR、GDPval 与 DeepScholar-Bench 分别暴露时间跨度、经济输出和研究综合的不同瓶颈。}{01:02:10--01:04:08}

\videofigure{figures/benchmark-realwork-gap.jpg}{三类 benchmark 都与真实工作存在系统差异：自动评分、一次性任务、上下文缺失或来源约束。}{01:04:08--01:05:42}

\begin{table}[H]
\centering
\small
\begin{tabular}{>{\raggedright\arraybackslash}p{2.6cm}>{\raggedright\arraybackslash}p{3.4cm}>{\raggedright\arraybackslash}p{3.8cm}>{\raggedright\arraybackslash}p{3.5cm}}
\toprule
\textbf{评估} & \textbf{核心问题} & \textbf{主要优势} & \textbf{主要盲点} \\
\midrule
METR & 能稳定完成多长任务？ & 人类时间统一标定、可追踪趋势 & 偏软件任务、质量与价值较弱 \\
GDPval & 输出是否胜过专家产物？ & 多模态真实工作、经济解释直接 & 专家偏好昂贵、上下文被简化 \\
DeepScholar & 能否完成可信研究综合？ & 同时测检索、综合与引用 & 领域有限、全面性仍难自动判定 \\
\bottomrule
\end{tabular}
\end{table}

\videofigure{figures/known-unknown.jpg}{我们较确定模型在短、清晰、可验证任务上快速进步，但对高上下文、开放环境和全面检索仍缺乏信心。}{01:05:42--01:07:18}

\videofigure{figures/evaluation-takeaways.jpg}{课程总结强调可靠性、错误恢复、任务特征和互补指标，而不是追逐单一排行榜。}{01:07:18--01:09:12}

\begin{importantbox}{评估应成为 agent 的训练接口}
细粒度失败标签可以直接指导改进：规划错误进入 planner 数据，工具错误进入 tool policy，遗漏来源进入 retrieval curriculum，引用不一致进入 verifier 训练。好的 benchmark 不只排名，还应定位下一轮自我改进信号。
\end{importantbox}

\subsection{开放问题}

\begin{itemize}
    \item 如何测量跨天、跨人的组织协作，而非单次沙箱任务？
    \item 如何把隐性上下文、权限、沟通和合规约束放入可重复评估？
    \item 如何评估 agent 在失败后恢复，而不只是最终成功或失败？
    \item 如何检测 reward hacking、benchmark contamination 和 evaluator bias？
    \item 如何把任务长度、质量、成本、能耗与风险合并为可决策的 Pareto frontier？
\end{itemize}

\subsection{本章小结}

没有任何单一 benchmark 能代表通用 agent。完整画像需要任务跨度、可靠性、输出质量、经济成本、上下文获取、来源覆盖和错误恢复等互补维度。

